\exercisegroup

\begin{enumerate}
\def\labelenumi{\arabic{enumi})}
\item
  证明可度量化空间的每个覆叠都是可度量化的。
\item
  证明局部紧且在无穷处可数的空间的连通覆叠是局部紧的且在无穷处可数的。
\item
  \begin{enumerate}
  \def\labelenumii{\alph{enumii})}
  \tightlist
  \item
    局部平凡纤维化的积是局部平凡纤维化，仅当其中除有限多个外都是平凡的。
  \end{enumerate}
\end{enumerate}

\begin{enumerate}
\def\labelenumi{\alph{enumi})}
\setcounter{enumi}{1}
\tightlist
\item
  满射覆叠的纤维积是覆叠，仅当其中除有限多个外都是同构。
\end{enumerate}

\begin{enumerate}
\def\labelenumi{\arabic{enumi})}
\setcounter{enumi}{3}
\tightlist
\item
  设 B 为子空间 \(\{1/n \mid n \in \mathbf{N}^*\} \cup \{0\}\) （在 \([0,1]\) 中），并设 \(E = B \times \mathbf{N} - \{(0,0)\}\) 。
\end{enumerate}

\begin{enumerate}
\def\labelenumi{\alph{enumi})}
\item
  证明 E 赋有映射 \(pr_{1}: E \to B\) 后是 B 的平凡覆叠。
\item
  证明典范单射 \(E \hookrightarrow B \times \mathbf{N}\) 是 B 的覆叠的态射，但并不使 E 成为 \(B \times \mathbf{N}\) 的覆叠。
\item
  构造两个 B-空间 \((\mathrm{E}_{1}, p_{1})\) 与 \((\mathrm{E}_{2}, p_{2})\) 及一个 B-态射 \(f: E_{1} \to E_{2}\) ，使得 \((\mathrm{E}_{1}, f)\) 与 \((\mathrm{E}_{1}, p_{1})\) 是覆叠，但映射 \(p_{2}\) 不是分离的。
\item
  证明形如 \((1/n,n)\) 的偶（当 n 取遍 \(\mathbf{N}^{*}\) ）所成的集 U 在 \(B \times \mathbf{N}\) 中既开又闭，但 \((\mathrm{U},\mathrm{pr}_{1}|\mathrm{U})\) 不是 B 的覆叠。
\end{enumerate}

\begin{enumerate}
\def\labelenumi{\arabic{enumi})}
\setcounter{enumi}{4}
\tightlist
\item
  设 B 为空间 \(\mathbf{N} \times ]0,1[\) 的亚历山德罗夫紧化，\(o\) 为其无穷远点（“夏威夷耳环”）。
\end{enumerate}

\begin{enumerate}
\def\labelenumi{\alph{enumi})}
\item
  证明空间 B 是连通的、局部连通的且局部道路连通的。
\item
  设 \((r_{n})_{n\in\mathbf{N}}\) 为趋于 0 的严格递减实数序列。对所有 \(n\in\mathbf{N}\) ，以 \(C_{n}\) 表示数平面上以 \((r_{n},0)\) 为心、以 \(r_{n}\) 为半径的圆；设 P 为诸 \(C_{n}\) 的并。证明 B 同胚于 P，且从 B 到 P 上的每个同胚都把点 o 映为平面的原点。
\item
  对 \(n \in \mathbf{N}\) ，记 \(P_{n} = C_{0} \cup \cdots \cup C_{n}\) ，并设 \(f_{n}\) 为从 P 到 \(P_{n}\) 的唯一映射，使得 \(f_{n}(x) = x\) （对 \(x \in P_{n}\) ），否则 \(f_{n}(x) = 0\) 。证明 \(f_{n}\) 连续。
\end{enumerate}

设 E 与 \(\mathrm{E}'\) 为 P 的覆叠。证明典范映射 \(\mathrm{Mor}_{\mathrm{P}}(\mathrm{E}',\mathrm{E})\to \varinjlim_{n}\mathrm{Mor}_{\mathrm{P}}(\mathrm{E}_{\mathrm{P}_n}',\mathrm{E}_{\mathrm{P}_n})\) 是双射。

设 E 为 P 的覆叠；证明存在 \(n \in \mathbf{N}\) ，使得 E 同构于 \(f_{m}^{*}(E_{\mathrm{P}_{m}})\) （对每个整数 \(m \geqslant n\) ）。

\begin{enumerate}
\def\labelenumi{\arabic{enumi})}
\setcounter{enumi}{5}
\tightlist
\item
  我们重新采用习题 5 的记号。
\end{enumerate}

\begin{enumerate}
\def\labelenumi{\alph{enumi})}
\item
  证明对每个整数 \(n \geqslant 1\) ，存在 P 的次数为 n 的覆叠 \((\mathrm{E}_{n}, p_{n})\) ，使得 \(p_{n}^{-1}(\mathrm{C}_{n})\) 是连通的。
\item
  设 E 为诸 \(E_{n}\) 的不交并，\(f: E \to P \times \mathbf{N}\) 为典范映射。证明 \((\mathrm{E}, f)\) 是 \(P \times \mathbf{N}\) 的覆叠，映射 \(pr_{1}\) 使 \(P \times \mathbf{N}\) 成为 P 的覆叠，但 \((\mathrm{E}, \mathrm{pr}_{1} \circ f)\) 不是 P 的覆叠。
\item
  设 \(P^{+}\) 与 \(P^{-}\) 分别为 P 中由点 \((x,y)\) （满足 \(y \geqslant 0\) 与 \(y \leqslant 0\) ）所成的集。证明 \(E|P^{+}\) 是 \(P^{+}\) 的平凡覆叠，\(E|P^{-}\) 是 \(P^{-}\) 的平凡覆叠。
\end{enumerate}

\begin{enumerate}
\def\labelenumi{\arabic{enumi})}
\setcounter{enumi}{6}
\tightlist
\item
  设 E、F、G 为拓扑空间，\(f: E \to F\) 与 \(g: F \to G\) 为连续映射，使得 \((E, f)\) 、\((F, g)\) 与 \((E, g \circ f)\) 是覆叠。
\end{enumerate}

\begin{enumerate}
\def\labelenumi{\alph{enumi})}
\item
  给出一个例子，其中 \(f\) 与 \(g \circ f\) 有次数，但 \(g\) 没有次数。
\item
  给出一个例子，其中 g 与 \(g \circ f\) 有次数，但 f 没有次数。
\end{enumerate}

\begin{enumerate}
\def\labelenumi{\arabic{enumi})}
\setcounter{enumi}{7}
\tightlist
\item
  \begin{enumerate}
  \def\labelenumii{(\arabic{enumii})}
  \setcounter{enumii}{1}
  \tightlist
  \item
    设 \(n\) 为自然数 \(\geqslant 1\) 。
  \end{enumerate}
\end{enumerate}

\begin{enumerate}
\def\labelenumi{\alph{enumi})}
\item
  对每个元素 \(x=(x_{1},\ldots,x_{n})\in\mathbf{C}^{n}\)，证明存在唯一的复系数首一多项式 P，次数为 \(n+1\)，使得 \(P(0)=0\) 且 \(x_{1},\ldots,x_{n}\) 是 P 的导数的零点。记 \(\theta(x)=(P(x_{1}),\ldots,P(x_{n}))\)；其坐标为 P 的临界值。
\item
  证明 \(\theta\) 定义了一个从 \(\mathbf{C}^n\) 到 \(\mathbf{C}^n\) 的多项式映射，满足 \(\theta^{-1}(0) = 0\) 。
\item
  以 U 表示 \(\mathbf{C}^{n}\) 的这样的开子集：其元素是坐标互异且非零的元素 \((x_{1},\ldots,x_{n})\) 。证明 \(\theta^{-1}\left(\mathrm{U}\right)\subset\mathrm{U}\) ，且 \(\theta\) 在 \(\theta^{-1}\left(\mathrm{U}\right)\) 上的限制定义了 U 的一个覆叠 \((\theta^{-1}\left(\mathrm{U}\right),\theta_{\mathrm{U}})\) 。
\item
  证明覆叠 ( \(\theta^{-1}(U),\theta_{U}\) ) 的次数为 \((n+1)^{n}\) 。
\item
  证明映射 \(\theta\) 是满射。
\end{enumerate}

\begin{enumerate}
\def\labelenumi{\arabic{enumi})}
\setcounter{enumi}{8}
\tightlist
\item
  设 B 与 F 为拓扑空间，E 为 B 上以 F 为纤维型的局部平凡纤维丛。
\end{enumerate}

\begin{enumerate}
\def\labelenumi{\alph{enumi})}
\item
  假设 B 是仿紧的，且 F 同胚于紧空间族的和空间。证明这时 E 是仿紧的。（重新采用 I, 第 84 页, 命题 8 的证明。）
\item
  假设 B 与 F 都是可度量化的。证明 E 也是可度量化的。
\item
  假设 B 与 F 都在无穷处可数。证明 E 在无穷处可数。
\end{enumerate}

